\section{Discussion}
\label{sec:discussion}

\subsection{Why the gap between disruption and injection is what it is}

We see three compounding causes:

\paragraph{Decoder fusion erases payload specifics.} Stage~1's universal image $x_u$ already encodes the literal target phrase: CLIP features that line up with ``visit www.example.com'' or ``card number''. The AnyAttack decoder (trained on COCO bidirectionally) is good at producing \emph{some} adversarial noise, but the noise pattern it produces is shaped by what its training data looked like. Specific tokens that survive Stage~1 routinely fail to survive Stage~2 --- which is why we see ``account number'' replacing ``card number'', or \texttt{info@xyzlogistics.com} replacing the chosen \texttt{support@fakecorp.com}.

\paragraph{The target \vlm{} needs a semantic invitation.} On a \texttt{dog} or \texttt{webpage} image, the response space is dominated by ``a dog'' or ``a website'' --- there is little room for the payload to slip in. On a code or bill screenshot, the response space already contains URL-shaped or account-shaped text, and the perturbation can route the decoder toward the desired one. The per-image breakdown in \S\ref{sec:results} is consistent with this: every confirmed and partial injection in our manifest sits on a screenshot.

\paragraph{Architecture wins over parameter count.} BLIP-2's Q-Former bottleneck filters out the precise low-level perturbations that drive Qwen-style \vlm{}s off-track. We did not retrain the AnyAttack decoder against BLIP-2, but the immunity is so complete (zero affected pairs out of 2,205) that we do not believe a budgeted, gradient-fused attack can puncture this architecture without enlarging the perceptual budget. This is, incidentally, why the v2 future work in \S\ref{sec:future} considers porting the Q-Former bottleneck as a defense layer.

\subsection{Why the v1 evaluation looked so much scarier}

Our first attempt at evaluation (v1, LLM-as-judge) reported a $50.5\%$ ``injection rate'' on Qwen2.5-VL-3B. The same data, scored with our v2 dual-dim evaluation, shows $100\%$ Output Affected and $0.41\%$ Target Injected. Where did the $50\%$ go? The judge prompt rewarded ``any deviation from the clean response'', so a Qwen response that called the dog photo a ``collage of text snippets'' counted as success even though the target phrase never appeared. The lesson is small but useful for the literature: if you only measure \emph{drift}, every successful attack on a small \vlm{} looks identical to a failed one --- which makes universal-attack methods look more dangerous than they are when used as a delivery vector.

\subsection{Cross-model transferability}

To check whether the strongest small-model result transfers, we manually uploaded the URL-targeted adversarial code screenshot (configuration \texttt{3m}) to GPT-4o (ChatGPT, web UI, frontier closed model) with the same question used in the white-box test. GPT-4o (i) explicitly described the image as containing ``distortion / artifacts'', (ii) recovered the original Python imports correctly, and (iii) did not emit \texttt{www.example.com}. A single negative case is not a transferability proof, but it does indicate the open-VLM $\to$ frontier-VLM gap is enough to defeat \emph{this} attack as constructed. Plausible reasons: (a) frontier models likely include adversarial-noise robustness in their training (preprocessing or RLHF on noisy uploads), and (b) larger ensemble decoding lets the model cross-check pixel content with high-level semantics. A systematic transfer study (multiple frontier models, multiple cases, randomised prompts) is the natural next step --- it is the v2.2 milestone in \S\ref{sec:future}.

\subsection{Limitations}

Three honest limits frame the scope of these claims. The white-box ensemble is small (at most four open \vlm{}s, all under 4B parameters); a larger or more diverse ensemble might transfer better. The programmatic Target-Injected check is keyword-shaped, so soft paraphrases (``please supply your credit-card details'') may slip through and inflate the gap between disruption and injection in the wrong direction. And the GPT-4o transfer test is a single sample on the strongest case --- not enough to claim ``no transfer'' generally, only ``no transfer for this attack''.

\subsection{The HuggingFace dataset as a project output}
\label{sec:discussion-dataset}

Beyond the report itself, the project ships a HuggingFace dataset (\url{huggingface.co/datasets/jeffliulab/visinject}) that bundles the $21$ universal images, the $147$ adversarial photos, the $6{,}615$ response pairs with v2 judge scores, and the $12$ curated injection examples. Figure~\ref{fig:hf-downloads} shows the download counter at $\sim 300$ in the first month. The dataset is small but specific: as far as we are aware, it is the first public release that bundles universal-attack outputs together with a dual-dimension evaluation, which makes it usable as a regression set for new \vlm{}s. Three plausible reasons it has been picked up: (i) the curated success cases save reviewers triage time; (ii) the pipeline reuses pretrained weights (\texttt{coco\_bi.pt}), so reproducing it is cheap; (iii) the dataset card explicitly documents the v1 vs v2 evaluation gap, which is a question other groups working on multimodal jailbreaks are likely to bump into.

\begin{figure}[h]
  \centering
  \includegraphics[width=0.65\textwidth]{HF-downloads.png}
  \caption{HuggingFace dataset download counter (screenshot taken April 2026, $\sim 300$ downloads). The download trace is the only piece of out-of-classroom feedback we have, and it suggests the artefacts have utility beyond the course report.}
  \label{fig:hf-downloads}
\end{figure}
